\documentclass[addpoints]{exam}
% \documentclass[addpoints,12pt]{exam}

\usepackage[slovene]{babel}
\usepackage{amsfonts}
\usepackage{amssymb}
\usepackage{amsmath}
\usepackage{float}
\usepackage{graphicx}
\usepackage{enumitem}
\usepackage{color}
\usepackage{eurosym}


%Komentar naslednje vrstice glede na verzijo testa -- rešitve ali prostor za odgovore:
% \printanswers

% \linespread{2}


\pointsinrightmargin
\marginbonuspointname{}


\renewcommand{\solutiontitle}{\noindent\textbf{Rešitve in točkovanje:}\par\noindent}

\colorgrids
\definecolor{GridColor}{gray}{0.7}
\setlength{\gridsize}{7mm}

\extrawidth{5mm}
\setlength{\rightpointsmargin}{1.5cm}

\begin{document}

\runningfooter{}{}{}

\begin{center}
    \fbox{\fbox{\parbox{15cm}{\centering
    Drugo pisno ocenjevanje znanja -- ponovno \\ 1. letnik -- test F \\ 19. december 2025 \\ (Potence in izrazi; Deljivost)}}}
\end{center}

\vspace{0.1in}
\makebox[\textwidth]{Ime in priimek, oddelek:\enspace\hrulefill}


\begin{center}
    \chqword{Naloga}
    \chpword{Možne točke}
    \chbpword{Dodatne točke}
    \chsword{Dosežene točke}
    \chtword{Skupaj}  
    % \combinedpointtable[h][questions]
    \combinedgradetable[h][questions]

\end{center}

\vspace{0.1in}


\begin{questions}

    % 1. vprašanje
    
    \question
    Izračunajte.
\begin{parts}
    \part[4]
    $\left(-1\right)^{2n+1}+7\cdot\left(-5+2\cdot 3\right)^{n+5}-\left(-1\right)^{2n+2}$

    \begin{solution}[4cm]
    $$\begin{aligned}
    \left(\left(-2\right)^3+\left(-5\right)^2\right)\cdot\left(-1\right)^{2025}-\left(\left(-3\right)^2\cdot\left(-2\right)-\left(-4\right)^2\right)-1^{2026}&=(-8+25)(-1)-(9(-2)-16)-1= \\ &=-17+34-1= \\ &=-18
    \end{aligned}$$
        Za pravilno uporabo pravila potenciranja z negativnim predznakom $1$ točka, pravilno potenciranje $1$ točka, pravilno množenje $1$ točka, končen rezultat $1$ točka.
    \end{solution}



    \part[4]
    $\left(xy^2z\right)^3\left(x^2y^3\left(-z^2\right)^3\right)\left(x^3\left(-y^2\right)^5(-z)\right)^3$

    \begin{solution}[4.5cm]
    $$5\cdot 2^{2n+1}-2\cdot 4^{n+1}-2\cdot 2^{2n+2}=5\cdot 2^{2n+1}-2^{2n+3}-2^{2n+3}=2^{2n+1}(5-2^2-2^2)=-3\cdot 2^{2n+1}$$
        Za pravilen zapis potenc števila $2$ $1$ točka, pravilno izpostavljanje skupnega faktorja $1$ točka, končen rezultat $1$ točka.
    \end{solution}

        
    \end{parts}


    % 2. vprašanje
    \question[5]
    Skrčite izraz in rezultat razstavite: $(x+1)^3-7x(x-1)+(x-4)^2-17$.

    \begin{solution}[5cm]
        \begin{align*}
            (7a-2b)^2-(5a+2b)(5a-2b)-6(4a^2-b^2)&=49a^2-28ab+4b^2-25a^2+4b^2-24a^2+6b^2=\\
                &=14b^2-28ab= \\
                &=14b(b-2a)
        \end{align*}
        Za pravilno kvadriranje binoma $1$ točka, za pravilno uporabo pravila razlike kvadratov $1$ točka, pravilno množenje $1$ točka, pravilno izpostavljanje $1$ točka, pravilno uporabljena Vietova formula $1$ točka.
    \end{solution}
    

    % \newpage
    % 3. vprašanje
    \question
    Faktorizirajte izraze, kolikor je mogoče.
    \begin{parts}
        \part[2]
        $x^2-4x-5$

        \begin{solution}[4cm]
            $$2x^2-30x+108=2(x-6)(x-9)$$
            Za uporabo Vietovega pravila $1$ točka, pravilen rezultat $1$ točka.
        \end{solution}


        \part[2]
        $y^3-16y$

        \begin{solution}[4cm]
            $$f^4-9g^2=(f^2-3g)(f^2+3g)$$
            Za pravilno izpostavljanje $1$ točka, pravilno uporabljena formula razlike kvadratov $1$ točka.
        \end{solution}


        \part[3]
        $4y^3-4y^2-9y+9$

        \begin{solution}[4cm]
            $$c^3-4c^2-c+4=(c-1)(c+1)(c-4)$$
            Za uporabo pravila faktorizacije štiričlenika 2+2 $1$ točka, pravilna faktorizacija štiričlenika $1$ točka, pravilno razstavljena razlika kvadratov $1$ točka.
        \end{solution}


        \part[3]
        $c^6+8c^3$

        \begin{solution}[4cm]
            $$m^5+27g^6m^2=m^2(m+3g^2)(m^2-6g^2+9g^4)$$
            Za pravilno izpostavljanje skupnega faktorja $1$ točka, uporaba formule razlike kubov $1$ točka, pravilen končen rezultat $1$ točka.
        \end{solution}


        \part[3]
        $x^5-32y^5$

        \begin{solution}[3cm]
            $$81-9b^2+12bc-4c^2=(9-3b+2c)(9+3b-2c)$$
            Za uporabo pravila faktorizacije štiričlenika 3+1 $1$ točka, pravilna uporaba formule $1$ točka, pravilno razstavljanje razlike kvadratov $1$ točka.
        \end{solution}


    \end{parts}


    %  \newpage

    % 4. vprašanje
    \question
    Določite neznano števko, da bo število $\overline{uuuuuuu4}$ deljivo
    \begin{parts}
        \part[2]
        s številom $4$.

        \begin{solution}[2cm]
            Za navedbo kriterija za deljivost z $2$ $1$ točka, uporaba kriterija in izračun $a\in\{0,2,4,6,8\}$ $1$ točka.
        \end{solution}


        \part[2]
        s številom $9$.

        \begin{solution}[4cm]
            Pravilno uporabljen kriterij za deljivost s $3$ $1$ točka, pravilno uporabljen kriterij za deljivost s $4$ $1$ točka, 
            pravilen zaključek in zapis kombinacij $(a,b)\in\left\{(2,0),(2,3),(2,6),(2,9),(6,2),(6,5),(6,8)\right\}$ $1$ točka.
        \end{solution}


    \end{parts}



% 5. vprašanje

    \question[4]
    Določite največji skupni delitelj $D$ in najmanjši skupni večkratnik $v$ izrazov $5x^2+20x-25$ in $25x^2-25$.

    \begin{solution}[6cm]
        Za zapisa $b=2\cdot 3^3\cdot 5^2\cdot 7$ in $c=2^2\cdot 3^2\cdot 5^4$ po $1$ točka. 
        Izračun $D(a,b,c)=3^2\cdot 5^2=225$ $1$ točka, in $v(a,b,c)=2^3\cdot 3^4\cdot 5^4\cdot 7\cdot 13$ $1$ točka.
    \end{solution}
    
    


    % 6. vprašanje
    \question
    Zapišite po besedilu in izračunajte število $x$ oziorma $y$.
    \begin{parts}
        \part[2]
        Če število $x$ delimo s število $235$, dobimo količnik $5$ in ostanek $25$.
        
    \begin{solution}[3cm]

    \end{solution}

        \part[3]
        Če število $5613$ delimo s število $y$, dobimo količnik $29$ in ostanek $16$.
    \end{parts}
    
    \begin{solution}[3cm]

    \end{solution}

        




     \newpage

    %dodatno vprašanje
    \bonusquestion[3]
    \textit{Dodatna naloga:} Dokažite: Vsota petih zaporednih naravnih števil je vedno deljiva s $5$.

    \begin{solution}[7cm]

    \end{solution}




\end{questions}



\end{document}